\documentclass[
  paperwidth=7in,
  paperheight=9in,
  blockwidth=6in,
  blockheight=8in,
  bleed=0.1in
]{photobook}
\usepackage[T1]{fontenc}
\usepackage{lmodern}
\usepackage{xcolor}
\usepackage{microtype}
\usepackage{hyperref}
\definecolor{night}{HTML}{17212B}
\definecolor{sunset}{HTML}{D97745}
\hypersetup{colorlinks=true,urlcolor=sunset,linkcolor=night}
\pagestyle{empty}
\title{A Small Atlas of\\Everyday Light}
\author{Your Name}
\date{2026}
\begin{document}
\maketitle
\chapter*{How to use this starter}
\addcontentsline{toc}{chapter}{How to use this starter}
This photo-book starter treats each page as a deliberate visual cell. Replace
the drawn panels with photographs using \verb|\imagecell[fill]{}{your-image}|,
or keep the panels as captions, maps, and small pieces of editorial text. The
\texttt{photobook} class supplies the page, cell, bleed, and spread mechanics;
the content below is intentionally synthetic.
\tableofcontents
\chapter{Morning routes}
\begin{inlinecell}{\textwidth}{2.1in}
  \colorbox{night}{\parbox[c][1.9in][c]{\textwidth}{\centering\color{white}
    \fontsize{25}{30}\selectfont Replace this panel\\[-2pt]
    \normalsize with a full-bleed photograph}}
\end{inlinecell}
\section*{06:40 — north window}
The first light is rarely a single colour. It moves across brick, glass, and
leaves at different speeds, which makes a sequence of small observations more
useful than one grand description. Use short sections like this for captions or
field notes alongside the visual material.
\begin{inlinecell}{0.47\textwidth}{1.35in}
  \colorbox{sunset}{\parbox[c][1.15in][c]{0.47\textwidth}{\centering\color{white}
    \Large IMAGE 01\\[-2pt]\small caption or location}}
\end{inlinecell}\hfill
\begin{inlinecell}{0.47\textwidth}{1.35in}
  \colorbox{night}{\parbox[c][1.15in][c]{0.47\textwidth}{\centering\color{white}
    \Large IMAGE 02\\[-2pt]\small caption or location}}
\end{inlinecell}
\chapter{After the weather}
\begin{inlinecell}{\textwidth}{2.4in}
  \fcolorbox{sunset}{white}{\parbox[c][2.1in][c]{0.94\textwidth}{\centering
    \color{night}\Large A place for a landscape,\\[0.3em]
    \normalsize contact sheet, or foldout reference}}
\end{inlinecell}
Pair a large visual cell with a short note. The class can align, crop, and add
bleed to real images while keeping the surrounding page structure stable. A
longer project can split each chapter into separate files and use the class's
spread helpers for print-ready layouts.
\chapter*{Colophon}
\addcontentsline{toc}{chapter}{Colophon}
Set the trim size, bleed, image resolution, and colour workflow for the printer
here. Keep a record of the photographer, location, date, and image licence for
every photograph you add to the book.
\end{document}
